\chapter{Valores de referência}


Os valores de referência constituem parâmetros técnicos utilizados como base comparativa para avaliar o desempenho energético da unidade consumidora. A partir dessa comparação, é possível estimar o potencial de economia de energia com a adoção de medidas de eficiência energética e classificar o nível de desempenho atual da edificação frente a padrões consolidados.

As referências adotadas neste relatório foram extraídas de fontes públicas e reconhecidas no setor elétrico, com foco em edificações públicas e institucionais. As principais são:

\begin{itemize}
    \item \textbf{Programa Nacional de Conservação de Energia Elétrica (Procel):} coordenado pela Eletrobras, reúne diretrizes, dados técnicos e experiências consolidadas em eficiência energética, com foco em edificações públicas e institucionais~\cite{procel_geral}.
    
    \item \textbf{Guia de eficiência energética em prédios públicos:} publicado pela Eletrobras/Procel em 2011, apresenta valores de referência de consumo específico por tipologia (escolas, hospitais, prédios administrativos), além de orientações práticas para a condução de diagnósticos e implementação de melhorias~\cite{procel_guia2011}.
    
    \item \textbf{Manuais do Programa de Eficiência Energética da ANEEL (PROPEE):} estabelecem os critérios técnicos e metodológicos para a elaboração de projetos no âmbito regulado, incluindo orientações para diagnóstico, estruturação de propostas e avaliação de resultados~\cite{propee_aneel}.
    
    \item \textbf{Relatórios de diagnósticos energéticos:} elaborados no contexto de chamadas públicas, como a CP Procel 001/2016 e a CP ANEEL 001/2017, reúnem dados reais de consumo e desempenho energético de unidades públicas em diferentes regiões do país~\cite{diagnosticos_chamadas}.
\end{itemize}

A Tabela~\ref{tab:valores-referencia} apresenta, a título de comparação, valores típicos de consumo específico por tipo de edificação, úteis para análise do desempenho energético da unidade em estudo.


\begin{table}[H]
\centering
\footnotesize
\renewcommand{\arraystretch}{1.3}
\begin{tabular}{|>{\columncolor[HTML]{1a3a5c}\color{white}}p{6cm}|p{6cm}|}
\hline
\textbf{Tipo de edificação} & \textbf{Consumo específico típico} \\ \hline
Escola pública (fundamental)              & 4 a 8 kWh/m²/mês \\ \hline
Unidade básica de saúde                   & 10 a 18 kWh/m²/mês \\ \hline
Hospital de médio porte                   & 250 a 400 kWh/leito/mês \\ \hline
Prédio administrativo                     & 8 a 12 kWh/m²/mês \\ \hline
Unidade esportiva / ginásio               & 3 a 7 kWh/m²/mês \\ \hline
Unidade com climatização central          & 15 a 22 kWh/m²/mês \\ \hline
\end{tabular}
\caption{Valores típicos de consumo específico por tipo de edificação pública}
\label{tab:valores-referencia}
\vspace{0.3em}
\begin{flushleft}
\footnotesize * Os valores são indicativos e podem variar conforme fatores climáticos, regime de operação, ocupação, tipo de equipamentos e práticas de manutenção.
\end{flushleft}
\end{table}

A comparação entre os indicadores reais da unidade e os valores de referência apresentados permite identificar desvios de desempenho, avaliando o grau de eficiência energética atual da edificação. Essa análise fornece uma base objetiva para a estimativa do potencial técnico de economia de energia, orientando a priorização das medidas propostas. Além disso, constitui um elemento fundamental para o dimensionamento das metas de economia a serem alcançadas com a implementação das ações, conforme estabelecido nos critérios metodológicos definidos no \textbf{Módulo 7 do Programa de Eficiência Energética da ANEEL (PROPEE)}~\cite{propee_aneel}, que tratam da quantificação de resultados e do cálculo da relação custo-benefício dos projetos.


% Exemplo de tabela de referência (caso queira inserir valores típicos):
% \begin{table}[H]
% \centering
% \renewcommand{\arraystretch}{1.3}
% \begin{tabular}{|>{\columncolor{azuluc}\color{white}}l|c|}
% \hline
% \textbf{Tipo de Edificação} & \textbf{Consumo Específico Típico (kWh/m²/mês)} \\ \hline
% Escolas públicas & 4,0 a 8,0 \\ \hline
% Unidades de saúde & 10,0 a 20,0 \\ \hline
% Prédios administrativos & 6,0 a 12,0 \\ \hline
% \end{tabular}
% \caption{Valores típicos de consumo específico por tipo de edificação}
% \label{tab:valores-referencia}
% \end{table}
